\documentclass{ip-lab}

\lablevel{n1}
\labnumber{Laboratorio 3}
\labtitle{Capas de lógica y presentación}

\makeheader

\begin{document}

\maketitle

\begin{sectionbox}{Actividad 1: Capa de lógica}
Escriba un módulo de Python con dos funciones:
\begin{enumerate}
    \item Una función que reciba un número de horas transcurridas desde el 1 de enero a las 00:00 y determine en qué fecha se estaría (mes y día), suponiendo que cada mes tiene 30 días.
    \item Una función que reciba un número de segundos transcurridos desde el 1 de enero a las 00:00 y determine en qué fecha se estaría (mes y día), suponiendo que cada mes tiene 30 días.
\end{enumerate}
Cada función debe retornar un \verb|str| con un formato de fecha legible para un transeúnte cualquiera.
\end{sectionbox}

\begin{sectionbox}{Actividad 2: Capa de presentación}
Para cada función de la capa de lógica, escriba un módulo que contenga una función que permita interactuar con el programa a través de la consola. Es decir:
\begin{enumerate}
    \item Cree un módulo \verb|consola_horas.py| que importe el módulo de lógica, solicite al usuario un número de horas y muestre un mensaje indicando la fecha correspondiente (mes y día) calculada desde el 1 de enero.
    \item Cree otro módulo similar que solicite al usuario un número de segundos y muestre un mensaje indicando la fecha correspondiente calculada desde el 1 de enero.
\end{enumerate}
\end{sectionbox}

\begin{sectionbox}{Actividad 3: Pruebas}
Ejecute cada consola para probar que el programa funciona. Puede usar los siguientes valores de prueba:
\begin{itemize}
  \item \verb|1_000| horas después del 1 de enero a las 00:00, la fecha es el 12 de febrero (a las 16:00 horas, pero eso no lo pide el enunciado).
  \item \verb|10_000_000| de segundos después del 1 de enero a las 00:00, la fecha es el 26/04.
  \item \verb|100_000_000| de segundos después del 1 de enero a las 00:00, la fecha es el 18/03 (aunque 3 años después, pero eso tampoco lo pide el enunciado).
\end{itemize}

\end{sectionbox}

\pagebreak

\begin{sectionbox}{Entrega}
  Cree un archivo comprimido .zip que contenga exactamente el módulo de la capa de lógica y los módulos de la capa de presentación. Entregue el archivo comprimido a través de Bloque Neón en el laboratorio del Nivel 1 designado como ``L3: Capas de lógica y presentación''.
\end{sectionbox}

\end{document}
